% ANCHOR: after_label def:c11-kernels
\begin{remark}[Physics notation: the composite line along $D$ and $P_m$]
The composite kernel is the line of the block after Section 9.2: the gauge propagator $\langle A_mA_n\rangle$, the classical scale transformation $(2x^n\partial_n+2)A_m$ of the gauge field inserted on it, and the line $\langle A^*_m\,c\rangle$, which in the gauge-fixed language is the ghost propagator joined to the antighost derivative $\partial_m\bar c$ that replaces the BRST source $A^*_m$ in Feynman gauge. Along a translation $P_m=\partial_m$ the inserted transformation is $b^n\partial_nA_m$. In momentum space the two propagators of the same momentum give the factor $e^{-2\varepsilon p^2}/p^4$ of the square kernel, so the line runs from a ghost $c$ to a gauge field $A_m$ and behaves as $|p|^{-3}$ at small momentum. It exists because the Feynman-gauge fixing and the conformal transformation of the gauge field do not commute, $(Q^{GF})^2\neq0$; along the R-symmetry rotations it vanishes, the ghost being an R-singlet. The sector parts of $S^{(1)}(\Delta)$ and $S^{(1)}(P_b)$ are the scale transformation and the translation of the three field types separately: the scalars $A_I$, the gauginos $\lambda_\alpha{}^i,\bar\lambda_{\dot\alpha i}$, and the gauge field with the ghost.
\end{remark}

% ANCHOR: after_label def:c11-sp
\begin{remark}[Physics notation: regulator insertions and the composite pair kernel]
The special edge is the difference between the regularised and the bare contraction of a field with its BRST source, $e^{-\varepsilon p^2}-1$ times the field--source delta function: the regulator insertion by which the regularised BRST variation fails to be a derivation (block after Theorem \ref{thm:c4-obstruction}). A special-edge composite cycle is a one-loop graph through two constituent fields of a composite operator with one dilatation or translation insertion on the ghost--gauge line and one regulator insertion. The regulator insertion lives at proper times $t\le\varepsilon$, so the graph is a pure short-distance contribution: powers $\varepsilon^{-m}$, a finite part and no $\log\varepsilon$. The finite part is the composite pair kernel, a finite local two-body contribution of the gauge-fixing sector to the one-loop action of the dilatation and of the translations on composite operators; with the proper-time cutoff $\varepsilon^{-1/2}$ playing the role of the momentum cutoff (block after Proposition \ref{prop:c10-bound}), its value depends on the shape of the cutoff and is a scheme-dependent number.
\end{remark}

% ANCHOR: after_label thm:c11-local
\begin{remark}[Physics notation: locality of the gauge-fixing contribution]
In a Feynman-gauge computation of the one-loop scale transformation of a gauge-invariant composite operator, the graphs with the dilatation insertion on an internal gauge line split into graphs through two constituent fields and self-energy-type graphs on one field. The latter contain the composite two-point function, nonlocal on $A_m$, $\lambda$, $\bar\lambda$ and $c$ through $\log p^2$ of the field momentum (block after Theorem \ref{thm:c9-quadratic}). For a BRST-invariant operator the Slavnov--Taylor identity cancels these nonlocal terms against the graphs through pairs, and the whole insertion becomes the sum over pairs of constituent fields of the finite parts of the graphs with one regulator insertion: a local two-body operator, of the same form as the one-loop dilatation operator acting on neighbouring letters [B03c, (2.12)]. The one-body part of the one-loop dilatation is then the field-strength renormalisation of the letters together with the finite counterterm of the dilatation, as in the block after Theorem \ref{thm:c10-D1}.
\end{remark}

% ANCHOR: after_label prop:c11-tc0
\begin{remark}[Physics notation: the finite renormalisations of the scale transformation]
$\hbar F_1^{(1)}(\Delta)$ is the finite local counterterm of the dilatation. Up to BRST-exact field redefinitions its free quadratic part is a combination of the classical scale transformations $\varphi\mapsto(2x^m\partial_m+2w_\varphi)\varphi$ of the three field types, with the coefficients $\theta$ for the scalars $A_I$, $\theta_\lambda$ for the gauginos and $\theta_V$ for the gauge field with the ghost. Such a term changes the one-loop dilatation of a letter by a multiple of its classical scale transformation, a multiple of $x^m\partial_m+w$, and leaves the anomalous dimensions unchanged. Without reflection invariance the list has one further entry, the rotation $\lambda_\alpha{}^i\mapsto e^{i\varphi}\lambda_\alpha{}^i$, $\bar\lambda_{\dot\alpha i}\mapsto e^{-i\varphi}\bar\lambda_{\dot\alpha i}$ of the free gaugino, which is odd under the reflection of the Euclidean theory (block after Theorem \ref{thm:c9-quadratic}).
\end{remark}

% ANCHOR: after_label thm:c11-arity2
\begin{remark}[Physics notation: the commutator of $D$ and $P_m$ at one loop]
With the physics dilatation $D=-E=-\tfrac12\Delta$ (key), the arity-two equation along $(\Delta,P_b)$ is the one-loop form of $[D,P_m]=P_m$ on composite operators, up to BRST-exact terms. In a scheme in which the translations act classically, it states that the part of the one-loop dilatation that is proportional to the position, $x^m$ times a one-loop translation, is BRST exact: the one-loop dilatation of an operator at the point $x$ is its anomalous-dimension operator, with no further term in $x^m\partial_m$. This condition fixes the coefficients $\theta$, $\theta_\lambda$, $\theta_V$ of the scale transformations in the finite counterterm, one for each field type, by requiring that the multiple of $x^m\partial_m$ produced by the gauge-fixing graphs cancel.
\end{remark}

% ANCHOR: after_label thm:c11-families
\begin{remark}[Physics notation: the light-cone families and the universal density]
$n^m$ with $n^2=0$ is a light-cone direction, $D_n=n^mD_m$ is the light-cone covariant derivative $D_+$, $\psi=\lambda_1{}^i$ is the $+$ component of a Weyl gaugino, and $F_{11}$ is the self-dual field-strength component $F_{++}$; the letters are the covariant letters $D_+^kZ$, $D_+^k\psi$, $D_+^kF_{++}$ of the block after Definition \ref{def:c10-letters}. The families are the light-cone $\mathfrak{sl}(2)$ sectors of [B03c]: the twist-two scalar operators $\operatorname{tr}(D_+^kZ\,D_+^{S-k}Z)$, the gaugino operators $\operatorname{tr}(D_+^k\psi\,D_+^{S-k}\psi)$ and the gluon operators $\operatorname{tr}(D_+^kF_{++}\,D_+^{S-k}F_{++})$, whose one-loop anomalous dimensions are the harmonic sums $\gamma=\frac{g^2N}{2\pi^2}h(j)$ with $j=S$, $S+1$, $S+2$ [B03c, (2.31)], together with the $\mathfrak{su}(1|1)$ sector of $Z$ and one gaugino component and two mixed families. The theorem says that, on two neighbouring fields of a single trace, the gauge-fixing contribution acts as $-\tfrac43$ units times the identity on the same ordering, with no exchange and no trace term, plus $-\tfrac16$ units times the derivative of the pair along the translation: a universal multiple of the classical scaling, $-\tfrac23$ units for each field and each derivative.
\end{remark}

% ANCHOR: after_label cor:c11-D1fam
\begin{remark}[Physics notation: anomalous dimensions and the absent one-loop translation]
The one-loop dilatation is $\tfrac12D_1$ (block after Definition \ref{def:c8-components}), so on the families $\tfrac12D_1=\gamma$ with $\gamma=\frac{g^2N}{8\pi^2}E$ and $E$ an eigenvalue of Beisert's one-loop Hamiltonian $\sum_lH_{l,l+1}$, $H_{12}=2h(J_{12})$: $\gamma=\frac{g^2N}{2\pi^2}h(S)$ on twist two, $\gamma=\frac{3g^2N}{4\pi^2}$ on the Konishi multiplet, which contains the trace scalar bilinear, $\operatorname{tr}(F_{++}F_{++})$ and the $\mathfrak{su}(1|1)$ state $\operatorname{tr}(Z\psi\psi)$, and $\frac{g^2N}{8\pi^2}$ times the spectra of the $\mathfrak{su}(1|1)$ table. In the translation-covariant scheme $P_1=0$: the translations act on composite operators without a one-loop correction, $[P_m,O(x)]=\partial_mO(x)$, and the finite operator renormalisation of the block after Section 9.5 is not needed.
\end{remark}

% ANCHOR: after_label prop:c11-S
\begin{remark}[Physics notation: the special supercharges from the supercharges]
Classically the special supercharges $S^\alpha{}_i$, $\bar S^{\dot\alpha i}$ act on an operator at the point $x$ as the supercharges with the position-dependent parameter $\rho(x)\eta$, whose two-component form is $-i\,x_m\sigma^m_{\alpha\dot\alpha}\bar\eta^{\dot\alpha i}$ and $i\,x_m\eta^\alpha{}_i\sigma^m_{\alpha\dot\alpha}$, plus a term at the point: this is the content of the commutators $[P_m,S]\propto\sigma_m\bar Q$ and $[P_m,\bar S]\propto\sigma_mQ$ (Proposition \ref{prop:c2-special}). The proposition is the one-loop form of these commutators: when the one-loop translation vanishes, the position-dependent part of the one-loop special supercharge is the one-loop supercharge with the parameter $\rho(b)\eta$, and the one-loop special supercharge of an operator at the origin is the new datum. For the scalar bilinears, for which no such term at the point exists (block after Corollary \ref{cor:c7-S}), the one-loop $S$ is fixed by the one-loop $Q$. The one-loop $K_m$ is fixed by the one-loop $P_m$ and the one-loop normalisation of the two-point functions through $K=P^\dagger$ (Corollary \ref{cor:c8-gram}).
\end{remark}

% ANCHOR: after_label prop:c11-plane
\begin{remark}[Physics notation: the half-BPS bilinears of the twist]
For a rank-one supercharge $Q^i_\alpha u^\alpha z_i$-type, such as the supercharge $Q^4_-$ of [BK] and [CY13, (2.1)], the scalars with $\rho(w)\xi=0$ are the three complex scalars $A_{ij}z^j$, the letters $\gamma^I\sim\Phi^{4I}$ of [BK, (3.51)], annihilated by $Q$ classically. Their symmetric traceless bilinears $\operatorname{tr}(\gamma^{(I}\gamma^{J)})$ are half-BPS chiral primaries in the $\mathbf{20'}$ (block after Section 9.4). The proposition states that their one-loop $Q$-variation vanishes identically, at every order at which the selection rules of the block after Proposition \ref{prop:c7-select} hold, which is the untwisted form of $Q_1(\gamma^I\gamma^J)=0$ [BK, (3.77)].
\end{remark}

% ANCHOR: end_of_section 11.6
\begin{remark}[Physics notation: the twisted one-loop supercharge and the Konishi entry]
The twisted differential $Q^\xi_0+\hbar Q^\xi_1$ is the action of a single Poincar\'e supercharge $Q$ of rank one on composite operators, classically and at one loop, and its cohomology at $\hbar^0$ is the space of classically $\tfrac1{16}$-BPS operators [CY13]. The one-loop part acts on pairs of letters [BK, (3.7)] and raises the dimension by $\tfrac12$, a two-site operator of order $g^2N$ in the canonical normalisation with the leg count of the order-$g^2$ two-site terms of the $\mathfrak{su}(2|3)$ supercharge of [B03] (block after Section 10.9). The individual pair entries of [BK, (3.64)--(3.69)] depend on the normalisation of the operators: the Konishi entry, the one-loop term that maps the pair of a gaugino and a scalar to $\bar\lambda_{\dot\alpha}\bar\lambda^{\dot\alpha}$, is the twisted form of the one-loop Konishi anomaly, and its coefficient moves when the cubic scalar operator in the classical $Q$-variation is mixed with the operator $\operatorname{tr}(\bar\lambda_{\dot\alpha}\bar\lambda^{\dot\alpha})$ of the same quantum numbers. The scheme-independent object is the induced map $[Q^\xi_1]$ on the classical $Q$-cohomology (Theorem \ref{thm:c8-twist}), which decides which classically BPS operators acquire an anomalous dimension at one loop (block after Section 8.6).
\end{remark}
